\section{Related Work}
\label{sec:related}

\paragraph{Multi-agent collaboration and task performance.}
Frameworks such as AutoGen and MetaGPT organize LLM agents through conversation, specialized roles, and structured workflows~\cite{wu2023autogen,hong2024metagpt}. Multiagent debate studies how exchanges between model instances can improve factuality and reasoning~\cite{du2023improving}. These approaches make communication and division of labor central design choices. PreCast examines an additional consequence of those choices: how collaboration affects the authorization of executed actions. We distinguish approved-task completion from unauthorized execution because both can occur in the same episode. The motivating retail result concerns this joint outcome under a particular team architecture. It motivates predicting risk for a specified composition and protocol, with task utility evaluated alongside that prediction.

\paragraph{Interactive agent evaluation and failure analysis.}
AgentBench evaluates LLM agents across interactive environments, while $\tau$-bench studies tool--agent--user interaction in domains with task and policy requirements~\cite{liu2024agentbench,yao2024tau}. Such environments provide the states, tools, and observable outcomes needed to distinguish a successful authorized action from an unauthorized state change. MAST organizes failures observed in multi-agent trajectories into categories including system design, inter-agent misalignment, and task verification~\cite{cemri2026multi}. This retrospective perspective helps identify candidate behaviors for controlled probes. PreCast places an additional requirement on the resulting measurements: they must support a forecast fixed before the target team is evaluated. A failure taxonomy, a solo score, and a prospective team-risk forecast answer different evaluation questions. Our proposed comparison tests the added predictive value of controlled behavioral measurements over simpler summaries of models and teams.

\paragraph{Behavioral elicitation and sensitivity to social context.}
Model-written evaluations provide a way to elicit specific language-model behaviors systematically~\cite{perez2023discovering}. Work on sycophancy studies responses that align with a user's expressed views, making sensitivity to conversational context an explicit object of measurement~\cite{sharma2024towards}. PreCast applies controlled elicitation to operational decisions made in an assigned agent role. Its proposed probes vary the relationship between the original user's permission and a teammate's message: the teammate may convey permission accurately, make an unsupported approval claim, or present an earlier violation as a completed step. The observable response is whether the agent executes an unauthorized action, completes authorized work, or requests clarification. Transfer from these controlled contexts to naturally generated team interactions is an empirical question. In particular, behavioral interpretation alone does not establish that a measured response predicts the frequency of corresponding failures in an interacting team.

\paragraph{Prompt injection and conflicting instruction sources.}
AgentDojo evaluates prompt-injection attacks and defenses in tool-using agent environments~\cite{debenedetti2024agentdojo}. It provides a relevant perspective on how agents handle instructions encountered through sources with different authority. Delegation introduces a related distinction between a teammate's operational request and the user's permission to perform it. PreCast focuses on ordinary team interactions in which an unsupported instruction can arise from misunderstanding or faulty planning, without requiring a malicious participant. Conflicting peer messages appear in the controlled probes so that responses can be measured under comparable conditions. In the primary team evaluation, the collaboration itself generates the messages and any resulting conflicts. This separation allows us to test whether a susceptibility measured under controlled exposure forecasts violations that emerge during the workflow.

\paragraph{Runtime safeguards and prospective team selection.}
AgentSpec specifies and enforces runtime constraints on agent actions~\cite{wang2026agentspec}. SABER analyzes consequential state-changing steps and introduces safeguards including verification and targeted reflection around those steps~\cite{cuadron2025saber}. These approaches motivate strong intervention baselines for reducing unauthorized execution. PreCast's proposed decision point precedes the target team run: individual measurements are used to forecast risk and inform model selection, role assignment, or the choice of a safeguard. Runtime protection and prospective forecasting can be combined. The relevant evaluation is whether using a forecast leads to a better decision than a fixed composition or a uniformly applied safeguard, accounting for approved-task completion and cost. Establishing that benefit requires prospective selection and held-out execution; it cannot be inferred from the motivating team-versus-solo comparison alone.

